\section{Genetica mendeliana}

Gregor Mendel (1822--1884), il fondatore della genetica, condusse lo studio dei caratteri ereditari con esperimenti su piante di \textit{Pisum sativum}. I risultati che ottenne non solo furono straordinari, ma riuscì a interpretarli in maniera così corretta da fornire le basi a tutta la genetica moderna.

% ---------------------------------------------------------------------------
\subsection{Gli esperimenti di Mendel: incrocio monoibrido}

Incrocio riguardante l'altezza della pianta:
\begin{enumerate}
  \item varietà alte e nane di pisello vengono incrociate;
  \item la progenie ibrida (F$_1$) è tutta alta;
  \item la progenie ibrida F$_1$ viene autofecondata;
  \item nella progenie di questo incrocio (F$_2$) compaiono piante alte e piante nane in un rapporto approssimativo di 3:1 (787 piante alte, 277 piante nane).
\end{enumerate}

% ---------------------------------------------------------------------------
\subsection{I legge di Mendel o legge della dominanza}

Per ogni carattere considerato, nella F$_1$ derivata da un incrocio tra due diverse varietà, compariva soltanto un aspetto del carattere e mai l'altro. Il carattere che scompariva, o in qualche modo era nascosto nella F$_1$, riappariva nella F$_2$, ma soltanto con una frequenza pari a $\frac{1}{4}$ del numero totale di individui F$_2$.

Il carattere che si manifesta nella F$_1$ è detto \textbf{dominante} e prevale sull'altro, detto \textbf{recessivo}, che rimane latente nella F$_1$. Il carattere recessivo scompare apparentemente nella F$_1$, ma ricompare nella F$_2$ nella proporzione del 25\%: ciò avviene perché nella F$_1$ sono presenti entrambi i caratteri dei genitori, il dominante e il recessivo, e solo il dominante si esprime.

\rubrica{Rappresentazione simbolica dell'incrocio}
\begin{itemize}
  \item ciascun genitore omozigote produce un tipo di gameti: P, $DD$ (Alto) $\times$ $dd$ (Nano) $\rightarrow$ gameti $D$ e $d$;
  \item gli eterozigoti F$_1$ ($Dd$, Alto) producono due tipi di gameti in eguale proporzione;
  \item gli eterozigoti F$_1$ autofecondati producono una progenie F$_2$ alta e nana in rapporto 3:1: genotipi $DD$, $Dd$, $Dd$, $dd$ $\rightarrow$ rapporto genotipico $1\,DD : 2\,Dd : 1\,dd$, rapporto fenotipico $3$ alto $: 1$ nano.
\end{itemize}

% ---------------------------------------------------------------------------
\subsection{II legge di Mendel o legge della segregazione}

\rubrica{Prima ipotesi}
Le piante adulte contengono una coppia di ``fattori'' (geni) che determinano l'eredità di ciascun ``carattere'' (alleli).

\rubrica{Seconda ipotesi}
Se la coppia di geni presente in un certo individuo è costituita da due alleli diversi, uno dei due è dominante sull'altro.

\rubrica{Terza ipotesi}
Gli alleli di una coppia che determinano un carattere segregano (cioè si separano) al momento della formazione dei gameti $\rightarrow$ metà dei gameti contiene un allele, l'altra metà dei gameti contiene l'altro allele.

\rubrica{Base cromosomica della segregazione}
Il principio della segregazione di Mendel è correlato agli eventi della meiosi: la separazione dei cromosomi omologhi durante la meiosi ha come risultato la segregazione degli alleli in un eterozigote.
\begin{itemize}
  \item un gamete possiede un solo set di cromosomi, il numero $n$, e porta un cromosoma di ogni coppia di omologhi;
  \item perciò un determinato gamete può possedere solo un gene di ogni coppia di alleli;
  \item quando i gameti si uniscono, lo zigote che ne deriva è diploide ($2n$) e possiede coppie di cromosomi omologhi.
\end{itemize}

% ---------------------------------------------------------------------------
\subsection{Gli esperimenti di Mendel: incrocio diibrido}

Incrocio riguardante colore e forma del seme:
\begin{itemize}
  \item P: giallo, liscio $\times$ verde, rugoso;
  \item F$_1$ (autofecondata): giallo, liscio;
  \item F$_2$: giallo liscio, verde liscio, giallo rugoso, verde rugoso, in rapporto approssimativo $9:3:3:1$ (315, 108, 101, 32 individui).
\end{itemize}

\rubrica{Rappresentazione simbolica dell'incrocio}
\begin{itemize}
  \item ciascun genitore omozigote produce un tipo di gameti: P, $GGWW$ (giallo, liscio) $\times$ $ggww$ (verde, rugoso) $\rightarrow$ gameti $GW$ e $gw$;
  \item gli eterozigoti F$_1$ ($GgWw$, giallo liscio) producono quattro tipi di gameti in eguale proporzione: $GW$, $Gw$, $gW$, $gw$;
  \item gli eterozigoti F$_1$ autofecondati producono una progenie F$_2$ con quattro fenotipi in rapporto $9:3:3:1$.
\end{itemize}

\begin{table}[htbp]
  \centering
  \small
  \renewcommand{\arraystretch}{1.3}
  \begin{tabularx}{\textwidth}{|X|c|c|}
    \hline
    \textbf{Fenotipo} & \textbf{Rapporto genotipico} & \textbf{Rapporto fenotipico} \\
    \hline
    giallo, liscio & $GGWW{:}GGWw{:}GgWW{:}GgWw = 1{:}2{:}2{:}4$ & 9 \\
    \hline
    giallo, rugoso & $GGww{:}Ggww = 1{:}2$ & 3 \\
    \hline
    verde, liscio & $ggWW{:}ggWw = 1{:}2$ & 3 \\
    \hline
    verde, rugoso & $ggww$ & 1 \\
    \hline
  \end{tabularx}
\end{table}

% ---------------------------------------------------------------------------
\subsection{III legge di Mendel o legge dell'indipendenza}

Gli alleli che determinano due caratteri diversi segregano indipendentemente l'uno dall'altro al momento della produzione dei gameti.

Il principio dell'assortimento indipendente è una conseguenza diretta degli eventi della meiosi: due diverse coppie di cromosomi omologhi possono allinearsi in metafase I secondo due disposizioni alternative, per cui una cellula meiotica produce in egual proporzione i quattro possibili tipi di gameti.

% ---------------------------------------------------------------------------
\subsection{Dominanza incompleta}

Nella dominanza incompleta, gli alleli di un gene non sono né completamente dominanti né completamente recessivi. Il fenotipo dell'organismo eterozigote è diverso da quello di entrambi gli omozigoti.

\begin{itemize}
  \item fenotipo rosso, genotipo $WW$;
  \item fenotipo rosa, genotipo $Ww$;
  \item fenotipo bianco, genotipo $ww$.
\end{itemize}

% ---------------------------------------------------------------------------
\subsection{Codominanza}

Esempio: il sistema dei gruppi sanguigni ABO. Gli alleli $I^A$ e $I^B$ sono codominanti tra loro, ed entrambi dominanti sull'allele $i$.

\begin{table}[htbp]
  \centering
  \small
  \renewcommand{\arraystretch}{1.3}
  \begin{tabularx}{\textwidth}{|X|c|c|c|c|}
    \hline
    \textbf{Genotipo} & \textbf{Gruppo sanguigno} & \textbf{Antigene A} & \textbf{Antigene B} & \textbf{Frequenza (\%)} \\
    \hline
    $I^AI^A$ o $I^Ai$ & A & $+$ & $-$ & 41 \\
    \hline
    $I^BI^B$ o $I^Bi$ & B & $-$ & $+$ & 11 \\
    \hline
    $I^AI^B$ & AB & $+$ & $+$ & 4 \\
    \hline
    $ii$ & O & $-$ & $-$ & 44 \\
    \hline
  \end{tabularx}
\end{table}

% ---------------------------------------------------------------------------
\subsection{Epistasi}

È il fenomeno per cui l'espressione di un gene dipende dall'effetto di un altro gene. Nell'interazione genica o \textbf{epistasi} (\textit{epi} $\rightarrow$ sopra; \textit{stasis} $\rightarrow$ stare) i geni interagiscono: uno o più alleli di un gene inibiscono o mascherano l'effetto di uno o più alleli di un altro gene $\rightarrow$ alcuni fenotipi attesi non si manifestano nella discendenza.

Il gene che maschera l'espressione di un altro gene è detto \textbf{epistatico}; quello la cui espressione viene mascherata è detto \textbf{ipostatico}.

\rubrica{Esempio: colore del pelo nel Labrador}
\begin{itemize}
  \item un gene codifica per un enzima coinvolto nella \textit{produzione} di melanina: allele dominante $B$ ($BB$ o $Bb$) $\rightarrow$ colore del pelo nero; allele recessivo $b$ ($bb$) $\rightarrow$ colore del pelo marrone;
  \item un altro gene controlla la \textit{deposizione} del pigmento nel pelo: allele dominante $E$ ($EE$ o $Ee$) $\rightarrow$ pelo nero nei cani $BB$ o $Bb$, pelo marrone nei cani $bb$; allele recessivo $e$ ($ee$) $\rightarrow$ pelo giallo indipendentemente dal genotipo di $B$;
  \item incrocio tra genitori omozigoti $BBEE \times bbee$ $\rightarrow$ cuccioli F$_1$: $BbEe$.
\end{itemize}

\begin{table}[htbp]
  \centering
  \small
  \renewcommand{\arraystretch}{1.3}
  \begin{tabularx}{\textwidth}{|X|X|c|}
    \hline
    \textbf{Genotipo} & \textbf{Fenotipo} & \textbf{Frequenza in F$_2$} \\
    \hline
    $B\_E\_$ & pelo nero & $9/16$ \\
    \hline
    $bbE\_$ & pelo cioccolato & $3/16$ \\
    \hline
    $\_\_ee$ & pelo giallo & $4/16$ \\
    \hline
  \end{tabularx}
\end{table}
